% Baccalauréat malgache, série C, 2009 : énoncé.
\TCDBacAttribution{C}{2009}
\providecommand{\dx}{\mathrm{d}x}
\providecommand{\Re}[1]{\mathrm{Re}\left( #1 \right)}
\providecommand{\Im}[1]{\mathrm{Im}\left( #1 \right)}
\providecommand{\dt}{\mathrm{d}t}
\providecommand{\Ouv}{(O;\vec{u},\vec{v})}
\providecommand{\mes}[2]{\mathrm{mes}\left( \widehat{ \overrightarrow{#1} ; \overrightarrow{#2} } \right)}
\providecommand{\D}{(\mathscr{D})}
\providecommand{\Courbe}{(\mathscr{C})}
\providecommand{\congru}[3]{#1 \equi #2 \quad \left[ \mathrm{mod}~#3\right]}
%
%%%%%%%%%%%%%%%%%%%%%%%%%%%%%%%%%%%%%%%%%%%%%%%%%%%%%%%%%%%%%%%%%%%%%%%%%%%%%%%%%%%%%%%%%%%%%%%%%%%%%%%%%%%%%%%%%%%%%%%%%%%%%%%%
\noindent \underline{\textsc{\textbf{Exercice}}} (4 points) \\%[-1\baselineskip]
      \begin{description}
            \item[A- Probabilit\'e :]
                  Chaque lettre du mot « FLEURS » est écrite sur chaque face d’un dé cubique normal 
                  (une seule lettre sur chaque face). 
                  Un enfant lance trois fois successive ce dé et d’une manière independante. 
                  Il obtient ainsi un « mot » de trois lettres (un mot peut avoir un sens ou non).
                  \begin{enumerate}
                        \item Calculer la probabilité de chacun des événements suivants :
                              \begin{itemize}
                                    \item $A$ : « L’enfant obtient un mot commençant par une consonne ».
                                          \pnt{0,5}
                                    \item $B$ : «  L’enfant obtient un mot commençant et se terminant par la même lettre ».
                                          \pnt{0,5}
                              \end{itemize}
                        \item Soit $C$ l’événement « l’enfant obtient un mot ayant des lettres toutes différentes ».
                              \begin{enumerate}
                                    \item Définir l’événement $\overline{C}$, événement contraire de $C$.
                                          \pnt{0,5}
                                    \item Calculer la probabilité de l’événement $\overline{C}$.
                                          \pnt{0,5}
                              \end{enumerate}
                  \end{enumerate}
                  %
                  \begin{center}
                        NB: On donnera les résultats sous forme de fraction irréductible.
                  \end{center}
            \item[B- Arithm\'etique :]
                  \begin{enumerate}
                        \item Dans cette question $a$ désigne un élément de $\mathbb{Z}/6\mathbb{Z}$.
                              Calculer $a^2 + 3a + 2$ pour $a = 4$. 
                              \pnt{0,25}
                        \item \begin{enumerate}
                                    \item R\'esoudre dans $\mathbb{Z}^2$ l'\'equation $(E) : 7x - 3y = 0$.
                                    \item D\'eterminer, en utilisant l'algorithme d'Euclide, une solution particuli\`ere
                                          de l'\'equation 
                                          \begin{equation}
                                          \TCDtag{F}
                                                7 x - 3 y = 2 \;.
                                          \label{eq: Bac C 2009 Arithm F}
                                          \end{equation}
                                          R\'esoudre dans $\mathbb{Z}^2$, l'\'equation \eqref{eq: Bac C 2009 Arithm F}.
                              \end{enumerate}
                  \end{enumerate}
      \end{description}
%%%%%%%%%%%%%%%%%%%%%%%%%%%%%%%%%%%%%%%%%%%%%%%%%%%%%%%%%%%%%%%%%%%%%%%%%%%%%%%%%%%%%%%%%%%%%%%%%%%%%%%%%%%%%%%%%%%%%%%%%%%%%%%%
\noindent \underline{\textsc{\textbf{PROBLEME I}}} (7 points) \newline
      Dans le plan orient\'e $\Plan$, on donne le parall\'elogramme $OABC$.
      
            \begin{tikzpicture}[line cap=round,line join=round,>=triangle 45,x=1.0cm,y=1.0cm]
                  \clip(-1.59,-0.54) rectangle (6.45,2.54);
                  \draw (-1,0)-- (4,0);
                  \draw (4,0)-- (6,2);
                  \draw (6,2)-- (1,2);
                  \draw (1,2)-- (-1,0);
                  \begin{scriptsize}
                        \fill [color=black] (-1,0) circle (1.5pt) node[above] {$A$};
                        \fill [color=black] (4,0) circle (1.5pt) node[above] {$B$};
                        \fill [color=black] (6,2) circle (1.5pt) node[below] {$C$};
                        \fill [color=black] (1,2) circle (1.5pt) node[below] {$O$};
                  \end{scriptsize}
            \end{tikzpicture}
            
      \underline{\textbf{Partie A}} \newline
      Construire les triangles $OCD$ et $OEA$, rectangle et isoc\`eles en $O$, de sens direct.
      Puis placer le point $G$, milieu du segment $[DE]$.
      
      \underline{\textbf{Partie B}} \newline
      On se propose de d\'emontrer par deux m\'ethodes diff\'erentes que les vecteurs 
      $\overrightarrow{CA}$ et $\overrightarrow{OG}$ sont orthogonaux et que $CA = 2 OG$.
      \begin{enumerate}
            \item \underline{Par les nombres complexes.} \newline
                  On munit le plan $\Plan$ d'un rep\`ere orthogonal direct d'origine $O$.
                  On d\'esigne par $a$ et $b$ les affixes respectives des points $C$ et $E$.
                  \begin{enumerate}
                        \item D\'eterminer les affixes respectives des points $D$ et $A$.
                        \item D\'eterminer $Z$ et $Z'$ affixes respectives des vecteurs 
                              $\overrightarrow{CA}$ et $\overrightarrow{OG}$.
                        \item Exprimer $Z'$ en fonction de $Z$.
                        \item D\'emontrer que $\overrightarrow{CA}$ et $\overrightarrow{OG}$ sont orthogonaux et que $CA = 2 OG$.
                  \end{enumerate}
            \item \underline{Par les transformations.} \newline
                  On donne : 
                  \begin{mitanisa}
                        $\bullet\ r$ la rotation de centre $O$ et d'angle $\frac{\pi}{2}$ ; \\
                        $\bullet\ h$ l'homoth\'etite de centre $D$ et de rapport $2$ ; \\
                        $\bullet\ s$ la compos\'ee $r \circ h$.
                  \end{mitanisa}
                  \begin{enumerate}
                        \item Donner la nature de la transformation $s$.
                        \item D\'eterminer $s(G)$ et $s(O)$.
                        \item En d\'eduire que les vecteurs $\overrightarrow{CA}$ et $\overrightarrow{OG}$
                              sont orthogonaux et que $CA = 2 OG$.
                  \end{enumerate}
      \end{enumerate}
      
      \underline{\textbf{Partie C}} \newline
      Dans le rep\`ere orthogonal $(O ; \overrightarrow{OC} , \overrightarrow{OD})$, d\'eterminer l'expression complexe
      de la transformation $s$ ainsi que l'affixe de son centre.
%%%%%%%%%%%%%%%%%%%%%%%%%%%%%%%%%%%%%%%%%%%%%%%%%%%%%%%%%%%%%%%%%%%%%%%%%%%%%%%%%%%%%%%%%%%%%%%%%%%%%%%%%%%%%%%%%%%%%%%%%%%%%%%%

%%%%%%%%%%%%%%%%%%%%%%%%%%%%%%%%%%%%%%%%%%%%%%%%%%%%%%%%%%%%%%%%%%%%%%%%%%%%%%%%%%%%%%%%%%%%%%%%%%%%%%%%%%%%%%%%%%%%%%%%%%%%%%%%
\noindent \underline{\textsc{\textbf{PROBLEME II}}} (9 points) \newline
      Soit $f$ la fonction d\'efinie sur $\mathbb{R}$ par $f(x) = x - 1 + (x^2 + 2) \expo{-x}$.
      On d\'esigne par $\Courbe$ sa courbe repr\'esentative dans un rep\`ere orthonorm\'e \Oij{} d'unit\'e $2\cm$.
      
      \underline{\textbf{Partie A}} 
      \begin{enumerate}
            \item On donne la fonction num\'erique $g$ d\'efinie sur $\mathbb{R}$ par $g(x) = 1 - (x^2 - 2x + 2) \expo{-x}$.
                  \begin{enumerate}
                        \item Dresser le tableau de variation de $g$.
                        \item D\'emontrer que l'\'equation $g(x) = 0$ admet une solution unique, 
                              not\'ee $\alpha$ dans $\mathbb{R}$ et que $0,35 \leq \alpha \leq 0,36$.
                        \item En d\'eduire le signe de $g(x)$ suivant les valeurs de $x$.
                  \end{enumerate}
            \item \begin{enumerate}
                        \item Montrer que pour tout $x$ de $\mathbb{R}$, on a : $f'(x) = g(x)$.
                        \item Dresser le tableau de variation de $f$.
                        \item D\'emontrer que la droite $\D$ d'\'equation $y=x-1$ est une asymptote de $\Courbe$.
                        \item Tracer dans le m\^eme rep\`ere, la droite $\D$ et la courbe $\Courbe$.
                  \end{enumerate}
                  (On prend, pour construire $\Courbe$, $\alpha = 0,35$).
      \end{enumerate}
      
      \underline{\textbf{Partie B}} \textit{Etude d'une suite.} \newline
      On appelle $I$ l'intervalle $[1;2]$.
      \begin{enumerate}
            \item On consid\`ere la fonction $h$ d\'efinie sur $I$ par : $h(x) = f(x) - x$.
                  \begin{enumerate}
                        \item D\'emontrer que $h$ est strictement d\'ecroissante sur $I$,
                              puis dresser son tableau de variation.
                        \item En d\'eduire que l'\'equation $f(x) = x$ admet une solution unique, 
                              not\'ee $\beta$, appartenant \`a $I$.
                  \end{enumerate}
            \item D\'emontrer que, pour tout r\'eel $x$ de $I$, on a : $f(x) \in I$ et $|f'(x)| \leq \frac{3}{4}$.
            \item $(U_n)_{n\in\mathbb{N}}$ est la suite num\'erique d\'efinie par : $U_0 = 1$ et $U_{n+1} = f(U_n)$ o\`u 
                  $n\in\mathbb{N}$.
                  \begin{enumerate}
                        \item D\'emontrer par r\'ecurrence que, pour tout entier naturel $n$, on a : $U_n \in I$.
                        \item D\'emontrer que, pour tout entier naturel $n$, on a : 
                              $
                              |U_{n+1} - \beta| \leq \frac{3}{4} |U_n - \beta|
                              $
                              et que 
                              $
                              |U_n - \beta| \leq \left( \frac{3}{4} \right)^n |U_0 - \beta|
                              $.
                        \item Prouver que la suite $(U_n)$ converge vers un r\'eel $\ell$, que l'on pr\'ecisera.
                  \end{enumerate}
            \item D\'eterminer l'entier naturel $p$ tel que, pour tout entier naturel $n$ sup\'erieur ou \'egal \`a $p$, on ait :
                  $|U_n - \beta| \leq 10^{-2}$.
      \end{enumerate}
%%%%%%%%%%%%%%%%%%%%%%%%%%%%%%%%%%%%%%%%%%%%%%%%%%%%%%%%%%%%%%%%%%%%%%%%%%%%%%%%%%%%%%%%%%%%%%%%%%%%%%%%%%%%%%%%%%%%%%%%%%%%%%%%



































